% !TeX root = ../tesis.tex


\chapter{Esferocito}

\label{section:apendix4}

En la Sección \ref{section:Res_Esferocito} se discutieron los resultados obtenidos de $Q_{\text{ext}}$ y del patrón de esparcimiento para $\lambda=431$ nm, correspondiente a la banda de Soret de un esferocito de radio 2680~nm y una CHCM de 34 g/dL. En esta sección se presentan los resultados para las longitudes de onda representativas de las bandas $\beta$ y $\alpha$. En la Fig.~\ref{fig:FF_541_esf_1} se muestra la sección transversal diferencial de esparcimiento $\dd{C_{\text{sca}}}/\dd{\Omega}$, así como su representación polar [Fig. \ref{subfig:FF_polar_541_esf}] para los cortes $\varphi=0^{\circ}$ y $\varphi=90^{\circ}$ para $\lambda=541$ nm. En la Fig. \ref{fig:FF_cartesian_541_esf} se muestra la representación cartesiana de $\dd{C_{\text{sca}}}/\dd{\Omega}$ para los cortes $\varphi=0^{\circ}$ y $\varphi=90^{\circ}$ para $\lambda=541$ nm. Al igual que para $\lambda=431$~nm, el esparcimiento se concentra predominantemente en la región frontal, alcanzando un valor máximo de $g=0.99$. Por otra parte, la diferencia entre los cortes $\varphi=0^{\circ}$ y $\varphi=90^{\circ}$ se localiza principalmente en la región lateral, donde la correlación aumenta a $r=0.90$. 
	%
\begin{figure}[h!]
	\centering
	\sidesubfloat[First image]{\hspace{-0.5cm}{	\includegraphics[width=7.3cm]{4-Resultados/FFdistribution541_Esf.pdf}\label{subfig:FF_541_esf}}}\hspace{0.1cm}
	\sidesubfloat[Second image]{\hspace{-0.5cm}{
			\includegraphics[width=8cm]{4-Resultados/polarplot541FFComEsfEri.pdf}\label{subfig:FF_polar_541_esf}}}\vspace{-0.1cm}
	\caption{\textbf{a)} Sección transversal diferencial de esparcimiento $\dd{C_{\text{sca}}}/\dd{\Omega}$ normalizada respecto a su valor máximo para una longitud de onda $\lambda=541$ nm en función del ángulo de esparcimiento $\theta$ y del ángulo $\varphi$ para un esferocito. Se identifican dos planos $\varphi=0^{\circ}$ en rojo claro y $\varphi=90^{\circ}$ en rojo oscuro. \textbf{b)} Sección transversal diferencial de esparcimiento $\dd{C_{\text{sca}}}/\dd{\Omega}$ en decibeles en función de $\theta$ para un esferocito (en rojo) y un eritrocito sano (azul) para dos cortes: $\varphi=0^{\circ}$ (puntos color claro) y $\varphi=90^{\circ}$ (puntos color oscuro). }
	\label{fig:FF_541_esf_1}
\end{figure}
%

	%
\begin{figure}[h!]
	\centering
			\includegraphics[width=14cm]{4-Resultados/linearplot541FFEsf.pdf}\vspace{-0.1cm}
	\caption{Sección transversal diferencial de esparcimiento $\dd{C_{\text{sca}}}/\dd{\Omega}$ en función de $\theta$ para un esferocito (en rojo) y un eritrocito sano (azul) para dos cortes: $\varphi=0^{\circ}$ (puntos color claro) y $\varphi=90^{\circ}$ (puntos color oscuro).}
	\label{fig:FF_cartesian_541_esf}
\end{figure}
%

En la Fig.~\ref{subfig:FF_573_esf} se muestra la sección transversal diferencial de esparcimiento $\dd{C_{\text{sca}}}/\dd{\Omega}$, así como su representación polar [Fig. \ref{subfig:FF_polar_573_esf}] para los cortes $\varphi=0^{\circ}$ y $\varphi=90^{\circ}$ para $\lambda=573$~nm. En la Fig. \ref{fig:FF_cartesian_573_esf} se muestra la representación cartesiana de $\dd{C_{\text{sca}}}/\dd{\Omega}$ para los cortes $\varphi=0^{\circ}$ y $\varphi=90^{\circ}$ para $\lambda=573$ nm. Al igual que para las longitudes de onda de 431 y 541 nm, el esparcimiento se concentra predominantemente en la región frontal con $g=0.99$. Sin embargo, las principales diferencias entre los cortes $\varphi=0^{\circ}$ y $\varphi=90^{\circ}$ se localizan en la región de retroesparcimiento, donde se tiene que $r=0.79$, siendo el menor valor obtenido para las tres longitudes de onda. 
	%
\begin{figure}[h!]
	\centering
	\sidesubfloat[First image]{\hspace{-0.5cm}{	\includegraphics[width=7.3cm]{4-Resultados/FFdistribution573_Esf.pdf}\label{subfig:FF_573_esf}}}\hspace{0.1cm}
	\sidesubfloat[Second image]{\hspace{-0.5cm}{
			\includegraphics[width=7.7cm]{4-Resultados/polarplot573FFComEsfEri.pdf}\label{subfig:FF_polar_573_esf}}}\vspace{-0.1cm}
	\
	\vspace{0.1cm}
	\caption{\textbf{a)} Sección transversal diferencial de esparcimiento $\dd{C_{\text{sca}}}/\dd{\Omega}$ normalizada respecto a su valor máximo para una longitud de onda $\lambda=573$ nm en función del ángulo de esparcimiento $\theta$ y del ángulo $\varphi$ para un esferocito. Se identifican dos planos $\varphi=0^{\circ}$ en rojo claro y $\varphi=90^{\circ}$ en rojo oscuro. \textbf{b)} Sección transversal diferencial de esparcimiento $\dd{C_{\text{sca}}}/\dd{\Omega}$ en decibeles en función de $\theta$ para un esferocito (en rojo) y un eritrocito sano (azul) para dos cortes: $\varphi=0^{\circ}$ (puntos color claro) y $\varphi=90^{\circ}$ (puntos color oscuro).}
	\label{fig:FF_573_esf_1}
\end{figure}
%
	%
\begin{figure}[h!]
	\centering
	
	\includegraphics[width=14cm]{4-Resultados/linearplot573FFEsf.pdf}\vspace{-0.1cm}
	\caption{Sección transversal diferencial de esparcimiento $\dd{C_{\text{sca}}}/\dd{\Omega}$ en función de $\theta$ para un esferocito (en rojo) y un eritrocito sano (azul) para dos cortes: $\varphi=0^{\circ}$ (puntos color claro) y $\varphi=90^{\circ}$ (puntos color oscuro).}
	\label{fig:FF_cartesian_573_esf}
\end{figure}
%

Se observa que las tres longitudes de onda muestran que el esparcimiento permanece predominantemente dirigido hacia adelante, como se refleja en el valor de $g$ entre 0.98 y 0.99 y en las contribuciones integradas de la región frontal resumidas en la Tabla~\ref{table:val_esf}. En contraste, la similitud entre los patrones laterales de ambos planos presenta un comportamiento no monótono con la longitud de onda, con un mínimo de correlación en $\lambda=431$ nm, mientras que los patrones en la región de retroesparcimiento muestran una disminución de $r$ conforme aumenta la longitud de onda.


